% Source: Weinberg GR, physical PDF pp. 378--379
%         (printed pp. 357--358).
% Coverage: Chapter 12 epigraph and opening overview before Section 12.1.
% Figures/tables/footnotes: chapter epigraph; no figures, tables, or footnotes.
% Status: source-reviewed and compile-clean.
% Uncertainties: none.

\begin{flushleft}
\begin{minipage}{0.37\textwidth}
\raggedright
\itshape
``You could validly argue that the minimum formulation is neat, but really no
better than the other formulation. However move from this lecture room to your
bathtub and observe your big toe in the water. Your limbs no longer appear
straight because the velocity of light in the water differs from that in air.
The least-time principle tells you how to formulate behaviour under such
conditions and the memorizing of Snell's law about angles does not. Who can
doubt which is the better scientific explanation?''

\medskip
\upshape
\textit{Paul E.~Samuelson, Maximum Principles in Analytical Economics, Nobel
Memorial Lecture, December 11, 1970}
\end{minipage}
\end{flushleft}

There are a great many physical systems whose dynamic equations can be
derived from a ``principle of least action,'' that is, from a statement that
some functional of the dynamical variables, the ``action,'' is stationary
with respect to small variations of these variables. This formulation of the
dynamic equations has one great advantage: It allows us to establish an
immediate connection between symmetry principles and conservation laws.

The symmetry of the action that concerns us most in this book is general
covariance. In this section we shall develop a general definition of the
energy-momentum tensor for any material system, as a functional derivative of
the action for that system. The use of the action principle and general
covariance will then allow us to show that this tensor is indeed conserved.

To achieve a truly general formulation of general relativity in terms of an
action principle, it is necessary to uncover a question that has been
carefully buried until now: How can we incorporate the effects of gravitation
into the field theories of particles with half-integer spin? The answer
requires the development of an approach to general relativity, the ``tetrad
formalism,'' based directly on the families of locally inertial frames that
were our starting point in Chapter~3. Although the proof is more complicated,
the energy-momentum tensor in this formalism is still symmetric and
conserved.
